\documentclass[a4paper,11pt]{article}
\usepackage[margin=2.5cm]{geometry}
\usepackage{amsmath,amssymb,amsfonts}
\usepackage{graphicx}
\usepackage{hyperref}
\usepackage{booktabs}
\usepackage{xcolor}

\hypersetup{colorlinks=true,linkcolor=blue!50!black,citecolor=blue!50!black,urlcolor=blue!50!black}

\title{\textbf{Numerical Audit of a Phenomenological Curvature-Bound Bounce Ansatz:\\
Stability Gates, Exclusions, and Observational Requirements}}

\author{Ho Hyung Kim\\[4pt]
\small\textit{Independent Researcher, Seoul, Republic of Korea}\\
\small\texttt{careye\char`\@khu.ac.kr; iconarea@gmail.com}}

\date{August 2026}

\begin{document}
\maketitle

\begin{abstract}
Classical general relativity can lead to singularities of divergent curvature and density in black-hole interiors and in a backward extrapolation of the early Universe. We audit a phenomenological framework consisting of a coordinate-invariant curvature bound and an effective non-singular bounce at a critical density. The constant-equation-of-state background coincides exactly with a background ansatz used in a published stable $c_T=1$ beyond-Horndeski construction, with an effective Friedmann residual of $2.8\times10^{-17}$. This reproduces a known class of backgrounds; it does not derive a CRBC-specific microscopic action. A phenomenological $w(t)$ extension supplies coefficient trajectories that pass the stated ghost, gradient, and subluminality gates, but its covariant action, cutoff, and non-linear stability remain open. The audit excludes the minimal single-field realization and the constant-equation-of-state baseline. The computed adiabatic spectrum has $n_s\simeq3.35$, incompatible with Planck; an added entropic sector can alter the tilt only with a tuned, externally supplied mass and free conversion efficiency. A perturbative Bianchi-I test shear produces a quadrupolar template with a measured scale dependence, while its amplitude remains free. No CMB map or observational likelihood is analysed. The outcome is therefore a reproducible set of exclusions and validation requirements, not a verified cosmological model.
\end{abstract}

\noindent\textbf{Keywords:} bounce cosmology; phenomenology; beyond-Horndeski; effective field theory; statistical anisotropy; numerical validation

%======================================================================
\section{Problem statement}
%======================================================================

In general relativity, gravity is spacetime curvature rather than a single force strength. The classical interior solution of a black hole and a backward extrapolation of FLRW cosmology both reach, under stated conditions, a singularity of divergent curvature. It is more conservative to read that divergence as a signal that the classical theory is insufficient in the region than as an actual infinity in nature.

One distinction matters throughout. The \textbf{event horizon} of a large black hole can be locally weakly curved and is not automatically a quantum-gravity region. The transition candidate considered here is not the horizon but the region where spacetime curvature or matter density approaches the Planck scale.

%======================================================================
\section{Assumptions and scales}
%======================================================================

\subsection{Planck-scale benchmark}

Combining $\hbar$, $G$ and $c$ gives the characteristic length and density
\begin{equation}
\ell_\mathrm{P}=\sqrt{\frac{\hbar G}{c^{3}}},
\qquad
\rho_\mathrm{P}=\frac{c^{5}}{\hbar G^{2}} .
\end{equation}
These are dimensional benchmarks at which quantum effects and gravity are expected to become simultaneously non-negligible, not a proof of a unified theory.

\subsection{Curvature-bound hypothesis}

A ``maximum gravity'' should be expressed through a coordinate-invariant curvature scalar rather than an observer-dependent acceleration. For the Kretschmann scalar $\mathcal K = R_{\mu\nu\rho\sigma}R^{\mu\nu\rho\sigma}$ we assume
\begin{equation}
\mathcal K \le \eta\,\ell_\mathrm{P}^{-4},
\label{eq:bound}
\end{equation}
with $\eta$ dimensionless. Equation~\eqref{eq:bound} is an \emph{assumption} of this framework and is not claimed to be derived from any known theory of quantum gravity.

\subsection{String length scale and the high-curvature crossover}

In perturbative string theory the fundamental length is $\ell_s=\sqrt{\alpha'}$. The low-energy derivative expansion is trustworthy for $\alpha'|R_{\mu\nu\rho\sigma}|\ll1$; when $\alpha'|R_{\mu\nu\rho\sigma}|\gtrsim1$, higher-curvature corrections and stringy states cannot be neglected. This motivates, but does not establish, the weaker hypothesis used here: that new degrees of freedom near the string scale may produce a finite effective curvature or a non-singular transition.

For observability, causality is imposed on the coherent interval rather than on the Universe as a whole,
\begin{equation}
L_{\mathrm{coh}}(t)\le D_{\mathrm{causal}}(t),
\end{equation}
where $D_{\mathrm{causal}}$ is the particle or event horizon of the chosen cosmology.

\subsection{Causality and unitarity}

The framework does not assume that causality is lost beyond an event horizon. It assumes instead that a more fundamental quantum causal rule or unitary evolution governs the transition region where the classical description fails. The explicit form of that rule is the central gap of the proposal.

%======================================================================
\section{An effective bounce ansatz}
%======================================================================

For a spatially flat homogeneous universe with $c=1$, the simplest phenomenological effective equation is
\begin{equation}
H^{2}=\frac{8\pi G}{3}\rho\left(1-\frac{\rho}{\rho_c}\right),
\qquad \rho_c=\xi\rho_\mathrm{P},
\label{eq:friedmann}
\end{equation}
with $\xi$ a dimensionless coefficient that a microscopic theory must supply. Combined with $\dot\rho+3H(\rho+p)=0$ this gives
\begin{equation}
\dot H=-4\pi G(\rho+p)\left(1-2\frac{\rho}{\rho_c}\right),
\end{equation}
so that for ordinary matter with $\rho+p>0$ one has $\dot H>0$ at $\rho=\rho_c$, permitting a non-singular transition from $H<0$ to $H>0$.

Equation~\eqref{eq:friedmann} is an \emph{illustrative ansatz} resembling effective forms widely studied in bouncing cosmology. Writing it down does not derive loop quantum gravity, string theory, or any other specific proposal. Section~\ref{sec:verification} reports what happens when one asks which microscopic realisations are actually compatible with it.

%======================================================================
\section{Black holes, the Big Bang, and a new expanding region}
\label{sec:bh}
%======================================================================

The minimal narrative of the hypothesis is: gravitational collapse or a black-hole interior evolves to high density; curvature and density approach the Planck bound; a quantum-gravity transition replaces the classical singularity; a bounce occurs; an expanding spacetime region follows.

If that final region is causally disconnected from the parent spacetime, an observer inside it could see a hot dense early universe resembling a Big Bang. This is \emph{not} the claim that the external Hawking evaporation of one black hole is observed as our Big Bang; the externally visible evaporation and the hypothesised interior bounce are distinct problems. Nor does the framework claim that a black hole promptly becomes a white hole in the same external universe: equation~\eqref{eq:friedmann} is a homogeneous ansatz and cannot simply be inserted into a black-hole interior. A separate spherically symmetric or rotating solution, a transition timescale, boundary conditions on energy and information, and an explicit causal structure would all be required.

Mass should likewise be described as a conversion between hot field and radiation energy and particle masses, not as creation from nothing. Which species are produced, and whether any stable electromagnetically inert remnant could constitute dark matter, requires a separate quantum-field-theoretic calculation.

%======================================================================
\section{Relation to other proposals}
%======================================================================

Conformal cyclic cosmology connects the far future of one aeon to the Big Bang of the next by a conformal rescaling. The model here does not assume a conformal matching; it focuses on a dynamical rebound at maximum curvature or critical density.

Ekpyrotic and cyclic models are closer, in that they treat a pre-Big-Bang contracting phase and a transition. This work does not assume brane collisions or a specific higher-dimensional theory. At this stage it is therefore not a replacement for those models but a statement of the singularity-avoidance and transition rules they all require.

%======================================================================
\section{Falsifiability and observational program}
\label{sec:falsify}
%======================================================================

To become a scientific model, the framework must make at least one quantitative prediction differing from standard inflationary cosmology.

\begin{enumerate}
\item \textbf{Primordial CMB fluctuations.} \emph{Partly supplied} (\S\ref{sec:verification}). The adiabatic tilt $n_s\simeq3.35$ is an exclusion; the quadrupolar modulation $g_*(k)\propto k^{-0.70}$ is a signal whose \emph{shape} is predicted. Its amplitude is not.
\item \textbf{Primordial gravitational waves.} \emph{Not supplied.} The surviving construction has $c_T^2=1$ by design and is therefore compatible with GW170817, but that is a constraint the construction was built to satisfy, not a prediction.
\item \textbf{Nucleosynthesis and structure formation.} A reheating history compatible with light-element abundances. Not addressed here.
\item \textbf{Black-hole observations.} A mass remnant, ringdown deviation, or echo signature cannot be claimed without an explicit matching rule, which \S\ref{sec:bh} does not supply.
\end{enumerate}

\subsection{Indirect JWST route}

JWST cannot observe Planck-scale curvature directly. It can constrain early black-hole seeds and their hosts --- for instance through a uniform census of $z\gtrsim4$ Little Red Dot candidates in public NIRCam and NIRSpec data, jointly fitting spectroscopic redshift, line width, host stellar mass and lensing selection. Such a program compares astrophysical seed channels; by itself it cannot establish a bounce, a wormhole, or a curvature bound.

%======================================================================
\section{Numerical verification}
\label{sec:verification}
%======================================================================

This section reports only those items of \S\ref{sec:falsify} that were actually computed. No observational data was used.

\subsection{The background lies inside a known stable family}

In units $\alpha=1$, $a_B=1$, the exact \emph{constant-$w$} solution of equation~\eqref{eq:friedmann} is
\begin{equation}
H(t)=\frac{t}{p\,(1+t^{2})},\qquad a(t)=(1+t^{2})^{1/(2p)},\qquad p=\tfrac32(1+w),
\end{equation}
which is \emph{exactly} the background ansatz of Ye and Piao \cite{YePiao} at constant $p$ and unit lapse --- a family for which they give an explicit $c_T=1$ beyond-Horndeski construction. The residual of equation~\eqref{eq:friedmann} on this solution is $2.8\times10^{-17}$: an identity, not a fit. This establishes only the constant-$w$ background correspondence. It does not derive the curvature bound, a CRBC-specific microscopic theory, or a black-hole-interior matching.

\subsection{Stability}

Two candidates are excluded. A minimally coupled single field $P(X)=KX+LX^{2}-V$ yields no point, out of $10^{6}$ sampled, that violates the null energy condition while remaining ghost-free and gradient-stable --- consistent with the Horndeski no-go theorem \cite{Kobayashi}. The constant-$w$ baseline fails for a subtler reason: scalar subluminality requires
\begin{equation}
c_s^{2}(+\infty)=p/3\le1 \iff w\le1,
\end{equation}
whereas suppressing anisotropy through contraction requires
\begin{equation}
\Delta\ln(\sigma^{2}/\rho)=(3/p-1)\ln(1+\eta_i^{2})<0 \iff w>1 .
\end{equation}
The two conditions do not overlap. Both boundaries are reproduced to four digits by an exhaustive parameter scan.

A phenomenological $w(t)$ transition profile ($w_i\to4.33$, $w_f=0.6$) supplies coefficient trajectories that pass the stated quadratic gates, as summarised in Table~\ref{tab:stability}. This is a necessary consistency check, not a proof that a covariant beyond-Horndeski/DHOST action produces the profile. Such an action, its background field equations, its cutoff, and its non-linear Bianchi-I stability must be established independently.

\begin{table}[t]
\centering
\begin{tabular}{lll}
\toprule
Condition & Value & Verdict \\
\midrule
No ghost (scalar) & $\min Q_s = 3.000$ & pass \\
No ghost (tensor) & $\min Q_T = 0.500$ & pass \\
Gradient stability & $\min c_s^{2} = 0.114$ & pass \\
Subluminality & $\max c_s^{2} = 0.849$ & pass \\
Tensor speed & $c_T^{2} = 1$ & by construction \\
Shear suppression & $\Delta\ln(\sigma^{2}/\rho) = -6.56$ & pass \\
EFT control & $\Lambda$ assumed, not derived & \textbf{not counted} \\
\bottomrule
\end{tabular}
\caption{Quadratic coefficient-gate results for the phenomenological $w(t)$ transition profile. They do not constitute a derivation from a covariant action. The final row passes only because a cutoff was supplied by hand and is not counted as a validation result.}
\label{tab:stability}
\end{table}

\subsection{String-scale cutoff check: a Planck-scale bounce is not controlled}

A revised numerical check replaces the arbitrary cutoff ratio with the explicit
identification $\Lambda=\rho_H^{1/4}$, where $\rho_H$ denotes a putative
Hagedorn/string density. For the direct identification $\rho_H=\rho_c$, the
same stable coefficient trajectory has
\begin{equation}
\max(E_{\rm char}/\Lambda)=0.519,
\end{equation}
and fails the declared $E_{\rm char}/\Lambda<0.1$ control gate. The failure is
not a ghost or gradient instability; it is a failure of the low-energy EFT to
control a bounce placed at the string scale.

The gate passes only after choosing $\rho_H/\rho_c=10^4$, for which
$\max(E_{\rm char}/\Lambda)=0.0519$. This is a useful negative result: the
calculation does not support a Planck- or string-scale controlled bounce. It
instead identifies a sub-string-scale phenomenological bounce as the regime in
which the chosen EFT coefficients remain controlled. No dynamics deriving the
ekpyrotic-to-Hagedorn transition, the density ratio, or the cutoff relation is
provided here.

\subsection{Primordial spectrum: the adiabatic mode is excluded}

Integrating $v_k''+(c_s^{2}k^{2}-z''/z)v_k=0$ with the $z''/z$ that follows from the surviving coefficient trajectory gives $n_s\simeq3.35$, incomparable with the Planck value $n_s=0.9649\pm0.0042$ \cite{Planck18}. This is a known property of ekpyrotic contraction rather than a surprise. The pipeline was validated against the analytic contracting-phase result
\begin{equation}
n_s-1=3-2\left|\beta-\tfrac12\right|,\qquad \beta=\frac{1}{p_i-1},
\end{equation}
which it reproduces to within $0.07$ over a factor five in $p_i$.

The adiabatic mode of the bounce therefore cannot originate the observed fluctuations. The standard remedy --- an entropic mechanism --- was tested on the same background. A tachyonic entropy mass $m_{\rm eff}^{2}=-\lambda H^{2}$ gives $n_s=0.9925$ at $\lambda=104$ and $n_s=0.9572$ at $\lambda=106.6$. Since $\mathrm{d}n_s/\mathrm{d}\lambda=-\beta^{2}/\nu$, matching Planck within its error bar fixes $\lambda$ to $0.29\%$. That tuning is the price of the repair, and no mechanism in this framework explains it.

\subsection{The coherent relic: shape predicted, amplitude not}

A classical relic obeys the same linear mode equation as the fluctuations, so its envelope is the initial profile multiplied by the transfer function, and only the transfer function is background dynamics. The measured transfer function carries \emph{no} band-limited feature: its cutoff lies two orders of magnitude below the comoving scale of the bounce, so every observable mode is far longer than the transition and carries no record of having crossed it. Splitting the transition into two timescales --- as a three-level synchronisation hierarchy would imply --- changes nothing, because producing a peak requires modes to exit, re-enter and exit again, and $|aH|$ remains monotonic in topology.

The angular structure is a stronger statement. On an isotropic FLRW background the mode equation depends on $|k|$ alone, so no direction dependence can be generated at all. Within this framework the only dynamical source of a direction is residual shear, and it \emph{forces}
\begin{equation}
P(k,\hat{\mathbf k})=P_{\rm iso}(k)\left[1+g_*(k)\,(\hat{\mathbf k}\cdot\hat{\mathbf n})^{2}\right],
\qquad
\frac{\mathrm{d}\ln g_*}{\mathrm{d}\ln k}=-0.70 .
\end{equation}
The multipole is fixed at $L=2$; the modulation is multiplicative rather than additive; and the tilt agrees across two shear amplitudes, so the shape is a prediction. The $\mu=(\hat{\mathbf k}\cdot\hat{\mathbf n})^{2}$ dependence was measured rather than assumed, and is linear to within $4\%$.

The free parameters of the observational template accordingly drop from six to two within this test setup: an overall size and an axis. The initial shear amplitude is \emph{not} fixed by the theory, so only an upper limit follows and detectability is not predicted. Using the Planck quadrupole constraint $\Delta g_*=0.016$ \cite{KimKomatsu}, the test-shear amplitude obeys $C\lesssim2.3\times10^{-6}$. That published limit assumes a scale-independent $g_*$, so this number is indicative pending a dedicated reanalysis with the predicted tilt and a full likelihood.

\subsection{What a black-hole-interior origin would require}

If the reading of \S\ref{sec:bh} is maintained, the level of anisotropy becomes the obstacle. A black-hole interior is Kantowski--Sachs rather than FLRW and starts at $\sigma^{2}/\rho\sim O(1)$. Reaching the regime verified above ($\sim10^{-7}$) demands $\Delta\ln(\sigma^{2}/\rho)\simeq-16$, which at $p_i=8$ corresponds to a contraction roughly two hundred times longer than the one computed here. The stability verdict would also have to be re-established non-perturbatively on a Bianchi~I background. The black-hole--Big-Bang reading is therefore neither excluded nor supported; what has changed is that the requirement is now quantified.

\subsection{What this section does not claim}

\begin{enumerate}
\item The effective-field-theory cutoff is not derived.
\item No microscopic mechanism fixes the relic amplitude; only its shape is predicted.
\item The entropy field and its $0.29\%$ tuning are inserted by hand.
\item No data analysis was carried out. The analysis design is pre-registered, but no CMB map was opened.
\end{enumerate}

%======================================================================
\section{Limitations and next steps}
%======================================================================

\begin{itemize}
\item \textbf{Open.} A microscopic derivation of $\eta$, $\xi$ and the transition dynamics. The EFT cutoff $\Lambda$ belongs here: it is the one entry of the stability gate that remains an assumption.
\item \textbf{Excluded in this EFT test.} Identifying the bounce density directly with a Hagedorn/string density: the resulting $E_{\rm char}/\Lambda\simeq0.519$ fails the stated control criterion.
\item \textbf{Open.} Covariant completion of the $w(t)$ profile. The present coefficient gates are necessary but do not substitute for deriving and checking the full background and perturbation equations of one action.
\item \textbf{Partly settled.} Stability including anisotropies. Quadratic scalar and tensor coefficients are positive for the supplied profile; anisotropy is computed only at the level of a perturbative test shear.
\item \textbf{Open.} Global matching from a black-hole interior to an expanding region.
\item \textbf{Open.} Information, entropy and Hawking radiation.
\item \textbf{Partly settled.} Observables comparable with data. The scalar tilt and the quadrupole shape are computed; the amplitude is not, and the gravitational-wave side is absent.
\end{itemize}

%======================================================================
\section*{Conclusion}
%======================================================================

A finite curvature or density scale remains a useful research hypothesis for replacing a classical singularity with a non-singular transition. The verification reported here moved the proposal in two directions at once. It advanced, in that the effective background belongs to a family with a stable microscopic realisation, and that a primordial spectrum and a quadrupole shape follow from it by calculation. It narrowed, in that the same calculation excluded the minimal single-field realisation, the constant-equation-of-state baseline, and the standing of the bounce's adiabatic mode as the origin of the CMB.

What survives is a test quadrupolar template whose shape follows from the supplied perturbative shear setup and whose size is not predicted. The value of this work lies not in declaring an answer to quantum gravity, but in recording which numerical checks pass, which candidate mechanisms are excluded, and which derivations and data analyses remain mandatory.

%======================================================================
\section*{Data and code availability}
%======================================================================

All numerical results were produced by scripts that read no observational data. The reproducibility package contains the code, gate criteria, simulation outputs, and pre-registered CMB analysis design. The public repository URL and archived release DOI are supplied with the release; no claim in this paper uses an unarchived observational fit.

\begin{thebibliography}{99}

\bibitem{YePiao}
G.~Ye and Y.-S.~Piao,
``Implication of GW170817 for cosmological bounces,''
\textit{Commun.\ Theor.\ Phys.} \textbf{71}, 427 (2019),
\href{https://arxiv.org/abs/1901.02202}{arXiv:1901.02202}.

\bibitem{Kobayashi}
T.~Kobayashi,
``Generic instabilities of non-singular cosmologies in Horndeski theory: a no-go theorem,''
\textit{Phys.\ Rev.\ D} \textbf{94}, 043511 (2016),
\href{https://arxiv.org/abs/1606.05831}{arXiv:1606.05831}.

\bibitem{Planck18}
Planck Collaboration,
``Planck 2018 results. VI. Cosmological parameters,''
\textit{Astron.\ Astrophys.} \textbf{641}, A6 (2020),
\href{https://arxiv.org/abs/1807.06209}{arXiv:1807.06209}.

\bibitem{KimKomatsu}
J.~Kim and E.~Komatsu,
``Limits on anisotropic inflation from the Planck data,''
\textit{Phys.\ Rev.\ D} \textbf{88}, 101301 (2013),
\href{https://arxiv.org/abs/1310.1605}{arXiv:1310.1605}.

\bibitem{PlanckIsotropy}
Planck Collaboration,
``Planck 2018 results. VII. Isotropy and statistics of the CMB,''
\textit{Astron.\ Astrophys.} \textbf{641}, A7 (2020),
\href{https://arxiv.org/abs/1906.02552}{arXiv:1906.02552}.

\bibitem{Battefeld}
D.~Battefeld and P.~Peter,
``A critical review of classical bouncing cosmologies,''
\textit{Phys.\ Rep.} \textbf{571}, 1 (2015),
\href{https://arxiv.org/abs/1406.2790}{arXiv:1406.2790}.

\bibitem{Novello}
M.~Novello and S.~E.~P.~Bergliaffa,
``Bouncing cosmologies,''
\textit{Phys.\ Rep.} \textbf{463}, 127 (2008),
\href{https://arxiv.org/abs/0802.1634}{arXiv:0802.1634}.

\bibitem{Penrose}
R.~Penrose,
``On the gravitization of quantum mechanics 2: conformal cyclic cosmology,''
\textit{Found.\ Phys.} \textbf{44}, 873 (2014).

\bibitem{McAllister}
L.~McAllister and E.~Silverstein,
``String cosmology: a review,''
\textit{Gen.\ Relativ.\ Gravit.} \textbf{40}, 565 (2008),
\href{https://arxiv.org/abs/0710.2951}{arXiv:0710.2951}.

\end{thebibliography}

\end{document}
